%=====================================================================
%  III. HASIL DAN PEMBAHASAN
%=====================================================================
\section{HASIL DAN PEMBAHASAN}

\subsection{Persamaan Klein-Gordon dalam Ruang-Waktu Schwarzschild}

Dari elemen garis pada Persamaan~\eqref{eq:metrik}, komponen tensor metrik invers
yang tak nol adalah $g^{tt} = -1/f(r)$, $g^{rr} = f(r)$, $g^{\theta\theta} = 1/r^{2}$,
dan $g^{\varphi\varphi} = 1/(r^{2}\sin^{2}\theta)$, sedangkan determinan metriknya
memberikan $\sqrt{-g} = r^{2}\sin\theta$. Substitusi besaran-besaran tersebut ke
dalam Persamaan~\eqref{eq:kg}, kemudian membagi hasilnya dengan $r^{2}\sin\theta$,
menghasilkan bentuk eksplisit persamaan Klein-Gordon,
\begin{align}
  -\frac{1}{f}\,\partial_{t}^{2}\phi
  &+ \frac{1}{r^{2}}\,\partial_{r}\!\left(f r^{2}\, \partial_{r}\phi\right)
   + \frac{1}{r^{2}\sin\theta}\,\partial_{\theta}
     \!\left(\sin\theta\, \partial_{\theta}\phi\right) \nonumber \\
  &+ \frac{1}{r^{2}\sin^{2}\theta}\,\partial_{\varphi}^{2}\phi = 0.
  \label{eq:kg-eksplisit}
\end{align}
Persamaan tersebut masih berupa persamaan diferensial parsial dengan empat variabel
bebas, sehingga belum dapat diselesaikan secara langsung.

\subsection{Separasi Variabel dan Persamaan Radial}

Karena latar belakang bersifat statis dan bersimetri bola, penyelesaian
Persamaan~\eqref{eq:kg-eksplisit} dapat dipisahkan dengan ansatz~\eqref{eq:ansatz}.
Bergantungan waktu yang bersifat harmonik memberikan faktor $-\omega^{2}/f$ pada
suku temporal, sedangkan suku sudut memenuhi persamaan eigen harmonik sferis,
\begin{equation}
  \frac{1}{\sin\theta}\,\partial_{\theta}
  \!\left(\sin\theta\, \partial_{\theta} Y_{\ell m}\right)
  + \frac{1}{\sin^{2}\theta}\,\partial_{\varphi}^{2} Y_{\ell m}
  = -\ell(\ell + 1)\, Y_{\ell m},
  \label{eq:eigen-harmonik}
\end{equation}
dengan $\ell(\ell+1)$ sebagai nilai eigennya \cite{riley2006}. Dengan demikian, seluruh
bergantungan sudut tereduksi menjadi sebuah konstanta pemisah, sehingga yang
tersisa hanyalah persamaan radial
\begin{equation}
  \frac{1}{r^{2}} \frac{d}{dr}\!\left(f r^{2} \frac{dR}{dr}\right)
  + \left(\frac{\omega^{2}}{f} - \frac{\ell(\ell+1)}{r^{2}}\right) R = 0.
  \label{eq:radial}
\end{equation}
Persamaan~\eqref{eq:radial} memperlihatkan bahwa persoalan empat dimensi telah
tereduksi menjadi persoalan satu dimensi dalam koordinat radial, yang merupakan
konsekuensi langsung dari simetri latar belakang.

\subsection{Persamaan Master dan Potensial Efektif}

Persamaan~\eqref{eq:radial} masih memuat suku turunan pertama dan koefisien yang
bergantung pada $r$, sehingga bentuknya belum sederhana. Kedua hal tersebut dapat
dihilangkan sekaligus melalui dua langkah. Pertama, fungsi radial diskalakan dengan
substitusi $R(r) = \Psi(r)/r$, yang menghasilkan
\begin{equation}
  \frac{f}{r}\frac{d^{2}\Psi}{dr^{2}} - \frac{f'}{r^{2}}\Psi
  + \left(\frac{\omega^{2}}{f} - \frac{\ell(\ell+1)}{r^{2}}\right)\frac{\Psi}{r} = 0.
  \label{eq:radial-Psi}
\end{equation}
Kedua, digunakan koordinat \textit{tortoise} pada Persamaan~\eqref{eq:tortoise-def},
sehingga operator turunan bertransformasi menurut $d/dr = f\, d/dr_{*}$. Dengan
mengalikan Persamaan~\eqref{eq:radial-Psi} dengan $r/f$ dan menerapkan transformasi
tersebut, diperoleh persamaan master
\begin{equation}
  \frac{d^{2}\Psi}{dr_{*}^{2}} + \left(\omega^{2} - V(r)\right)\Psi = 0,
  \label{eq:master}
\end{equation}
dengan potensial efektif
\begin{equation}
  V(r) = f(r)\left(\frac{\ell(\ell+1)}{r^{2}} + \frac{f'(r)}{r}\right).
  \label{eq:potensial}
\end{equation}
Persamaan~\eqref{eq:master} memiliki bentuk yang identik dengan persamaan
Schr\"{o}dinger tak gayut waktu, sehingga persoalan perturbasi lubang hitam dapat
diperlakukan sebagai persoalan hamburan gelombang satu dimensi oleh sebuah
potensial. Dengan menyubstitusikan fungsi metrik Schwarzschild pada
Persamaan~\eqref{eq:fungsi-metrik} beserta turunannya $f'(r) = 2M/r^{2}$, potensial
efektif memperoleh bentuk eksplisit
\begin{equation}
  V(r) = \left(1 - \frac{2M}{r}\right)
         \left(\frac{\ell(\ell+1)}{r^{2}} + \frac{2M}{r^{3}}\right).
  \label{eq:potensial-schw}
\end{equation}

Untuk kasus medan skalar bermassa, kehadiran suku massa pada
Persamaan~\eqref{eq:kg-bermassa} tidak mengubah simetri latar belakang, sehingga
prosedur separasi variabel dan transformasi yang sama tetap berlaku dan bentuk
persamaan master pada Persamaan~\eqref{eq:master} tidak berubah. Perbedaannya
muncul pada potensial efektif, yang memperoleh suku massa tambahan,
\begin{equation}
  V(r) = \left(1 - \frac{2M}{r}\right)
         \left(\frac{\ell(\ell+1)}{r^{2}} + \frac{2M}{r^{3}} + \mu^{2}\right).
  \label{eq:potensial-bermassa}
\end{equation}

\subsection{Karakteristik Potensial Efektif}

Potensial efektif pada Persamaan~\eqref{eq:potensial-schw} lenyap pada kedua ujung
domain fisis, yakni pada horizon peristiwa ($r \to 2M$, atau $r_{*} \to -\infty$)
karena faktor $f(r)$ lenyap di sana, dan pada jarak tak hingga
($r \to \infty$, atau $r_{*} \to +\infty$) karena kedua suku dalam kurung meluruh
sebagai $r^{-2}$ dan $r^{-3}$. Di antara kedua ujung tersebut potensial bernilai
positif, sehingga terbentuk sebuah penghalang potensial tunggal dengan satu puncak.
Bentuk inilah yang ditampilkan pada Gambar~\ref{fig:potensial-l}.

\begin{figure}[t]
  \centering
  \includegraphics[width=\columnwidth]{potensial-l.png}
  \caption{Potensial Efektif Medan Skalar Tak Bermassa terhadap Koordinat
  \textit{Tortoise} untuk Variasi Bilangan Momentum Sudut ($M = 1$)}
  \label{fig:potensial-l}
\end{figure}

Terlihat bahwa kenaikan bilangan momentum sudut $\ell$ meningkatkan ketinggian
puncak penghalang secara signifikan, sebagaimana dirangkum pada
Tabel~\ref{tab:puncak-l}. Hal ini merupakan konsekuensi langsung dari suku
$\ell(\ell+1)/r^{2}$ pada Persamaan~\eqref{eq:potensial-schw}, yang berperan sebagai
penghalang sentrifugal dan tumbuh secara kuadratik terhadap $\ell$. Secara fisis,
gangguan dengan momentum sudut yang lebih besar menghadapi penghalang yang lebih
tinggi, sehingga lebih sulit menembus daerah di sekitar lubang hitam. Letak puncak
juga bergeser sedikit menjauhi horizon dan mendekati $r = 3M$ seiring bertambahnya
$\ell$, yakni radius orbit lingkaran foton, karena pada $\ell$ besar suku
sentrifugal mendominasi suku kelengkungan $2M/r^{3}$.

\begin{table}[h]
  \centering
  \caption{Letak dan Ketinggian Puncak Potensial Efektif untuk Variasi $\ell$
  ($M = 1$, $\mu = 0$)}
  \label{tab:puncak-l}
  \fontsize{9}{11}\selectfont
  \begin{tabular}{ccc}
    \toprule
    $\boldsymbol{\ell}$ & $\boldsymbol{r_{\text{puncak}}}$
      & $\boldsymbol{V_{\text{maks}}}$ \\
    \midrule
    0 & 2,667 & 0,0264 \\
    1 & 2,886 & 0,0993 \\
    2 & 2,952 & 0,2471 \\
    3 & 2,974 & 0,4692 \\
    \bottomrule
  \end{tabular}
\end{table}

Pengaruh massa lubang hitam ditampilkan pada Gambar~\ref{fig:potensial-M}.
Bertambahnya $M$ menyebabkan penghalang potensial menjadi lebih rendah sekaligus
lebih lebar, dan letak puncaknya bergeser menjauh dari pusat. Perilaku ini dapat
dipahami dari sifat penskalaan Persamaan~\eqref{eq:potensial-schw}: potensial hanya
bergantung pada kombinasi $r/M$, sehingga $V \propto M^{-2}$ sedangkan letak
puncaknya $r_{\text{puncak}} \propto M$. Dengan kata lain, lubang hitam yang lebih
masif tidak menghadirkan penghalang dengan bentuk yang berbeda, melainkan
penghalang yang sama yang direntangkan pada skala panjang yang lebih besar dan
diturunkan ketinggiannya. Konsekuensinya, lubang hitam yang lebih masif memiliki
skala waktu karakteristik yang lebih panjang, sehingga gangguan yang terjadi
padanya berkembang lebih lambat.

\begin{figure}[t]
  \centering
  \includegraphics[width=\columnwidth]{potensial-M.png}
  \caption{Potensial Efektif Medan Skalar Tak Bermassa untuk Variasi Massa Lubang
  Hitam ($\ell = 2$)}
  \label{fig:potensial-M}
\end{figure}

Adapun pengaruh massa medan skalar ditampilkan pada Gambar~\ref{fig:potensial-mu}.
Berbeda dengan kedua parameter sebelumnya, kehadiran $\mu$ mengubah perilaku
asimtotik potensial. Pada Persamaan~\eqref{eq:potensial-bermassa}, ketika
$r \to \infty$ faktor $f(r)$ menuju satu dan kedua suku pertama dalam kurung
lenyap, sehingga potensial tidak lagi menuju nol melainkan menuju nilai konstan
$\mu^{2}$. Penghalang potensial dengan demikian tidak lagi berdiri di atas dasar
yang nol, melainkan di atas sebuah dataran (\textit{plateau}) setinggi $\mu^{2}$.
Secara fisis, hal ini menunjukkan bahwa medan skalar bermassa tidak dapat merambat
bebas menuju jarak yang sangat jauh sebagaimana medan tak bermassa; keberadaan
potensial yang tidak lenyap pada jarak jauh cenderung menahan sebagian medan tetap
berada di sekitar lubang hitam alih-alih meloloskannya sepenuhnya menuju tak
hingga.

\begin{figure}[t]
  \centering
  \includegraphics[width=\columnwidth]{potensial-mu.png}
  \caption{Potensial Efektif Medan Skalar Bermassa untuk Variasi $\mu$
  ($\ell = 1$, $M = 1$)}
  \label{fig:potensial-mu}
\end{figure}

Selain menaikkan dasar potensial, penambahan $\mu$ juga menaikkan puncak penghalang
dan menggesernya menjauhi horizon peristiwa. Perhitungan menunjukkan bahwa untuk
$\ell = 1$ dan $M = 1$, puncak penghalang terletak pada $r \approx 2{,}89$ ketika
$\mu = 0$, lalu bergeser ke $r \approx 3{,}21$ ketika $\mu = 0{,}3$, dan ke
$r \approx 3{,}67$ ketika $\mu = 0{,}4$. Bersamaan dengan itu, muncul lembah dangkal
di antara puncak penghalang dan dataran asimtotik. Ketika $\mu$ terus dinaikkan,
lembah tersebut semakin dangkal hingga akhirnya lenyap sama sekali pada
$\mu \approx 0{,}45$; di atas nilai tersebut potensial menjadi fungsi yang naik
secara monoton dari nol pada horizon peristiwa menuju $\mu^{2}$ pada jarak tak
hingga, sehingga struktur penghalang tunggal yang khas pada medan tak bermassa
tidak lagi terbentuk. Nilai ambang ini bergantung pada $\ell$, karena suku
sentrifugal turut menentukan seberapa menonjol puncak penghalang dibandingkan
dataran asimtotiknya.

\subsection{Pembahasan}

Hasil-hasil di atas memperlihatkan bahwa seluruh pengaruh geometri ruang-waktu
Schwarzschild terhadap dinamika medan skalar terangkum sepenuhnya di dalam satu
besaran, yaitu potensial efektif pada Persamaan~\eqref{eq:potensial-schw}. Struktur
persamaan master pada Persamaan~\eqref{eq:master} tidak bergantung pada latar
belakang mana pun yang ditinjau; yang berubah dari satu latar belakang ke latar
belakang lain hanyalah fungsi metrik $f(r)$ yang masuk ke dalam potensial. Sifat
ini menjelaskan mengapa perturbasi medan skalar dapat dipandang sebagai persoalan
hamburan gelombang satu dimensi yang standar, dan mengapa potensial efektif menjadi
objek fisis yang paling informatif dalam kajian ini.

Ketiga parameter yang ditinjau memberikan pengaruh yang secara kualitatif berbeda.
Bilangan momentum sudut $\ell$ mengubah ketinggian penghalang tanpa mengubah
perilaku asimtotiknya, massa lubang hitam $M$ mengubah skala panjang dan skala
ketinggian penghalang tanpa mengubah bentuknya, sedangkan massa medan skalar $\mu$
mengubah perilaku asimtotik potensial itu sendiri. Perbedaan ini penting karena
hanya $\mu$ yang mengubah watak persoalan secara struktural, yakni mengubah dasar
potensial pada jarak jauh dari nol menjadi $\mu^{2}$, sehingga medan tidak lagi
sepenuhnya bebas pada daerah asimtotik.

Perlu dicatat bahwa penurunan ini bertumpu pada dua asumsi pokok. Pertama, medan
skalar diperlakukan sebagai medan uji, sehingga pengaruh baliknya terhadap metrik
latar belakang diabaikan; asumsi ini hanya sahih selama amplitudo gangguan cukup
kecil. Kedua, latar belakang diasumsikan statis dan bersimetri bola, yang merupakan
syarat agar separasi variabel dapat dilakukan sebagaimana pada
Persamaan~\eqref{eq:eigen-harmonik}. Apabila salah satu asumsi tersebut dilanggar,
misalnya pada lubang hitam yang berotasi, persamaan gerak tidak lagi tereduksi
menjadi bentuk sesederhana Persamaan~\eqref{eq:master} dan diperlukan formalisme
yang lebih umum \cite{teukolsky1973}. Dalam kerangka asumsi tersebut, potensial
efektif yang diperoleh dalam penelitian ini konsisten dengan bentuk baku yang
dilaporkan dalam literatur perturbasi lubang hitam \cite{konoplya2011, berti2016}.
